%! AP Statistics — Unit 3, Lesson 3.5 — Warm-Up
\documentclass[10pt]{article}
\usepackage{apstats-article}
\usepackage{apstats-boxes}

\begin{document}

\pageheader{Unit 3, Lesson 3.5}{Warm-Up}
\namedateperiod

\vspace{0.3in}

% ── SECTION 1: Spiral Review ──────────────────────────────────────────────────
{\large\bfseries\color{navy} Part 1 \quad Spiral Review}

\vspace{0.12in}

\begin{enumerate}

    \item From Lesson 3.4 --- for each flawed survey, identify the type of bias and
    state whether the estimate is too high or too low.

    \vspace{8pt}
    \renewcommand{\arraystretch}{1.8}
    \begin{tabularx}{\linewidth}{@{} X p{3cm} p{2.2cm} @{}}
    \textbf{Survey} & \textbf{Bias type} & \textbf{Too high/low?} \\
    \hline
    A radio station asks listeners to call in if they support the new mayor. &
    \blank{2.8cm} & \blank{2cm} \\
    A phone survey uses only landline numbers. &
    \blank{2.8cm} & \blank{2cm} \\
    \end{tabularx}

    \vspace{0.18in}

    \item A researcher wants to know whether a new study app improves test scores.
    She lets students choose whether to use the app, then compares scores.
    \begin{enumerate}[label=(\alph*), itemsep=4pt]
        \item Is this an observational study or an experiment? \blank{4cm}
        \item What confounding variable might explain a difference in scores?\\
        \writeline
        \item What would the researcher need to change to establish that the app
        \emph{causes} higher scores?\\
        \writeline\writeline
    \end{enumerate}

\end{enumerate}

\vspace{0.25in}

% ── SECTION 2: Looking Ahead ──────────────────────────────────────────────────
{\large\bfseries\color{navy} Part 2 \quad Looking Ahead}

\vspace{0.12in}

\noindent Today we learn the components and principles of well-designed experiments.

\vspace{0.14in}

\begin{tcolorbox}[colback=greenbg, colframe=greenacc, boxrule=0.8pt, arc=2mm,
    left=4mm, right=4mm, top=3mm, bottom=3mm]
    \small
    \begin{itemize}[leftmargin=1.3em, itemsep=8pt]
  \item A well-designed experiment: comparison, random assignment, replication, control.
  \item Random \emph{assignment} (not just selection) is what allows causal conclusions.
  \item Blocking groups similar units together to reduce variability.
    \end{itemize}
\end{tcolorbox}

\vspace{0.2in}

\noindent\textcolor{slate}{\small Write one thing you already know about how experiments
work:}\\[8pt]
\writeline\\[2pt]\writeline

\end{document}
